% TOP_PROBLEM: {"id": 2579, "title": "Kourovka Notebook Problem 21.70", "category_id": 17, "category": {"id": 17, "name": "group_theory", "display_name": "Group Theory", "description": "Problems about groups, group actions, representations, and related algebraic structures.", "slug": "group-theory", "order_index": 17, "created_at": "2026-07-31T15:26:25.671Z"}, "status": "open", "problem_number": "KOU-21.70", "set_id": 10}
% FIRST_POSTED: 2026-10-01
% ENTRY_KIND: statement_only
% UnsolvedMath ID 2579; source catalogue code KOU-21.70
% !TeX program = lualatex
% Definition and sources only; this file is not an LLM solution attempt.
\documentclass[11pt,a4paper]{article}
\usepackage[margin=25mm]{geometry}
\usepackage{fontspec}
\setmainfont{lmroman10-regular.otf}[BoldFont=lmroman10-bold.otf,
 ItalicFont=lmroman10-italic.otf,BoldItalicFont=lmroman10-bolditalic.otf]
\usepackage{amsmath,amssymb,mathtools}
\usepackage{unicode-math}
\setmathfont{latinmodern-math.otf}
\AtBeginDocument{\let\setminus\smallsetminus}
\usepackage{xurl,hyperref}
\hypersetup{colorlinks=true,urlcolor=blue}
\setcounter{secnumdepth}{0}
\setlength{\parindent}{0pt}
\setlength{\parskip}{5pt}
\setlength{\emergencystretch}{3em}
\begin{document}
\section*{Kourovka Notebook Problem 21.70}\label{2579}
{\small UnsolvedMath ID 2579 · KOU-21.70 · Group Theory · Kourovka Notebook - New Problems, Issue 21}
\subsection{Definitions and mathematical statement}
For $R=\mathbb Z$ or a field, write $RG$ for the group ring and regard
$R$ as a trivial $RG$-module through the augmentation
$\sum r_g g\mapsto\sum r_g$. A group is of type $FP$ over $R$ here
if this module has a finite-length resolution by finitely generated
projective $RG$-modules. Projective means a direct summand of a free
module; the terms need not themselves be free.

Group cohomology with coefficients in the group ring is
$H^i(G;RG)=\operatorname{Ext}_{RG}^i(R,RG)$, computed by applying
$\operatorname{Hom}_{RG}(-,RG)$ to a projective resolution of $R$.
Right multiplication on the coefficient ring gives these cohomology
groups a right $RG$-module structure.
An orientable Poincar\'e duality group of dimension $n$ over $R$ is a
group of type $FP$ over $R$ satisfying
\[
 H^i(G;RG)=0\quad(i\ne n),\qquad H^n(G;RG)\cong R
\]
as right $RG$-modules, with trivial $G$-action on the right-hand $R$.
The module condition is the orientability requirement, not merely an
isomorphism of underlying abelian groups.

Fix $n\ge0$. If a group $G$ has this property in dimension $n$ over
every field, must it have the same property over $\mathbb Z$?
The dimension is common to all the field hypotheses; the resolutions
may depend on the field. The conclusion includes integral type $FP$
as well as integral cohomology vanishing and the trivial integral
dualizing module. Finite presentability of $G$ is not assumed.
\subsection{Short English statement}
Does orientable Poincaré duality in the same dimension over every field imply orientable integral Poincaré duality?
\subsection{Sources}
\begin{itemize}
\item[S1] ULAM.AI and the UnsolvedMath Contributors, supplied problems.json, record 2579 (KOU-21.70), Kourovka Notebook - New Problems, Issue 21. Catalogue identity and source transcription; CC BY 4.0.
\url{https://huggingface.co/datasets/ulamai/UnsolvedMath}.
\item[S2] The Kourovka Notebook, 21st edition, new problems, Problem 21.70. Expanded formulation of the supplied catalogue statement.
\url{https://arxiv.org/pdf/1401.0300v47}.
\end{itemize}
\end{document}
